\chapter{Narrative Analysis}
\label{ch:narr}

\section{Goal}

Detection accuracy alone does not explain \emph{what} misinformation is about.
For public-health communication, stakeholders need a map of recurring narratives---especially those linked to vaccine hesitancy and broader COVID distrust.
This chapter addresses \textbf{RQ3} using BERTopic on documents labeled as misinformation.

\section{Setup}

We aggregate misinformation texts from CONSTRAINT and MMCoVaR ($N=5{,}752$ after filtering).
Long news bodies are truncated to 800 characters for embedding efficiency.
Documents are embedded with Sentence-Transformers \texttt{all-MiniLM-L6-v2} and clustered with BERTopic (\texttt{min\_topic\_size}=15).
We obtain dozens of fine-grained topics; Table~\ref{tab:taxonomy} reports the top themes after qualitative renaming.

\section{Taxonomy of dominant narratives}

\begin{table}[H]
\centering
\caption{Top misinformation / hesitancy-related narratives (BERTopic + qualitative labels).}
\label{tab:taxonomy}
\small
\begin{tabular}{clrp{5.2cm}}
\toprule
\# & Theme & Size & Interpretation \\
\midrule
1 & Case/death sensationalism & 520 & Inflated or misleading severity narratives \\
2 & Politicized pandemic framing & 500 & Partisan blame / political weaponization \\
3 & Contested public measures (masks) & 214 & Contestation of NPIs mixed with rumor \\
4 & China / lab-origin claims & 188 & Origin conspiracies and geopolitical blame \\
5 & Miracle-cure / unproven drugs & 116 & HCQ and similar treatment myths \\
6 & Elite / Bill Gates vaccine distrust & 112 & Agenda and control conspiracies \\
7 & Hospital-overload visuals & 97 & Misleading crisis imagery \\
8 & Schools / children anxieties & 82 & Child-focused risk rumors \\
9 & Folk home remedies & 79 & Salt/steam/garlic ``cures'' \\
10 & Travel-ban rumors & 76 & False policy / mobility claims \\
\bottomrule
\end{tabular}
\end{table}

Approximately 2{,}000 documents fall into an outlier topic (too unique or noisy for stable clustering), which is expected for short social posts.

\section{Implications for vaccine hesitancy}

Although not every theme is narrowly about vaccine injection side effects, the taxonomy reveals a \emph{multi-narrative ecosystem} that can fuel hesitancy indirectly:
\begin{itemize}
  \item Distrust of institutions and elites (Theme 6) undermines confidence in vaccine recommendations.
  \item Miracle-cure narratives (Theme 5) reduce perceived need for vaccination.
  \item Sensational harm and hospital imagery (Themes 1, 7) amplify fear responses.
  \item Child/school themes (Theme 8) specifically target parental decision-making.
\end{itemize}
For communication campaigns, these clusters suggest differentiated messaging rather than a single ``correct the side-effect myth'' strategy.

\section{Link to the detection pipeline}

Operationally, the classifier can flag misinformation content at scale; BERTopic then summarizes \emph{what} is being flagged.
Together they form a monitoring loop useful for public-health analysts and platform trust-and-safety teams, even when multimodal fusion does not beat text on every dataset.
